\chapter{Introduction: Why Mathematics Is Structure}
\label{chap:I_introduction}

\setlength{\parindent}{0pt}

% =====================================================
% 1.1 Guiding Question
% =====================================================
\subsection{Guiding Question}
\index{mathematics}
\index{discovery vs. invention}
\index{natural sciences}
\index{Galilei, Galileo}

Why can the world be described mathematically at all?
Why do we need \textbf{numbers}, \textbf{forms}, and \textbf{functions}
in order to order and understand natural processes?
Is \textbf{mathematics} merely a tool invented by human beings,
or do we encounter in it a structure
that exists independently of us?

This question stands in the background of all natural science.
Already the Greeks discovered
that harmony in music and geometry is based on numerical relationships.
Later, \textbf{Galileo Galilei} formulated the idea
that nature is written in the language of mathematics.
And to this day it becomes evident again and again
that mathematical concepts,
which at first seem purely abstract,
capture real relationships with astonishing precision:
the paths of the planets,
the laws of motion,
the propagation of light,
or the structure of the quantum world.

Thus the guiding question touches on more than a method.
It leads directly to our image of reality:
Is mathematics a product of the human mind,
or is it the expression of an order
that we find and gradually uncover?

% =====================================================
% 1.2 Mathematics as a Universal Language
% =====================================================

\subsection{Mathematics as a Universal Language}
\label{sec:1.2_universal_language}
\index{universal language}
\index{cultures}
\index{zero}
\index{algebra}
\index{irrational numbers}

One remarkable feature of mathematics is that
its fundamental structures appear again and again,
independently of culture and historical period.
The Greeks discovered the \emph{irrational numbers},
Indian mathematicians introduced \emph{zero}
as a central element of calculation,
and Arabic scholars developed \emph{algebra}
into an independent method.
Although languages, symbols, and historical conditions differed,
the same inner order becomes visible in all these developments.

It is especially revealing
that different cultures often arrived at similar mathematical results
without direct exchange.
In number systems,
in geometric relationships,
or in computational methods for determining roots,
the same patterns appear again and again.
This speaks against understanding mathematics
as a mere product of arbitrary conventions.
Rather, it appears as a structure
that is not invented,
but discovered under different conditions.

In this sense, mathematics is a universal language.
It does not belong to any single culture,
because it is not bound to words,
but to relationships.
Its statements remain valid regardless of whether they are formulated
in Greek, Arabic, Sanskrit, or modern symbolic language.
Precisely this is the source of its special power:
it makes structures visible
that remain the same across times and cultures.

For this reason, mathematics remains the common language of the sciences.
It connects physics, computer science, engineering, and many other disciplines,
because it does not merely describe content,
but reveals the form of lawful relationships.
Where patterns are recognized,
formulated precisely,
and generalized,
mathematics emerges as a universal structure of knowledge.

\begin{NoteBox}[Independent Discoveries]
	\label{box:discovery}
	\small
	The fact that different cultures developed central mathematical ideas independently
	speaks against the notion of mathematics as a merely arbitrary invention.
	Rather, mathematics appears as a structure
	that is discovered again and again
	under different historical conditions.
\end{NoteBox}

% =====================================================
% 1.3 Experiments as a Key
% =====================================================

\subsection{Experiments as a Key}
\index{experiment}
\index{Galilei, Galileo}
\index{thermometer}
\index{linearity}
\index{proportionality}
\index{Ohm, Georg Simon}
\index{resistance}
\index{voltage}
\index{current}

Only through experiments does it become clear
which mathematical structure is contained in a natural phenomenon.
An experiment does not merely provide individual numerical values;
it reveals whether a lawful relationship exists between two quantities.

A particularly important case is \textbf{proportionality}.
If two quantities behave in such a way that doubling one quantity
also doubles the other,
then a proportional relationship is present.
In a diagram, this can be recognized by the fact
that the measurement points lie on a straight line through the origin.

This is exactly what \textbf{Georg Simon Ohm} found in his measurements.
He investigated how voltage $U$ and current $I$ relate to each other.
If the measured value pairs $(I,U)$ are plotted in a diagram,
then for an ohmic conductor the points lie on a straight line through the origin
(see Fig.~\ref{fig:ohm_proportionality}).
This shows:
the voltage $U$ is proportional to the current $I$.
Mathematically, one writes
\[
U \propto I.
\]

To turn this proportionality into an exact equation,
one introduces the constant factor of proportionality.
In Ohm's case, this factor is the \textbf{resistance} $R$.
This gives the relationship
\[
U = R \cdot I.
\]

Resistance is therefore not introduced arbitrarily,
but emerges from the experiment as the slope of the straight line in the diagram.
It describes how strongly a conductor impedes the flow of current.
The formula is thus not a mere invention,
but the mathematical condensation of a measured relationship.

\begin{figure}[h]
	\centering
	\begin{tikzpicture}[scale=1.0]
		
		% Axes
		\draw[->, thick] (0,0) -- (6,0) node[right] {Current $I$};
		\draw[->, thick] (0,0) -- (0,4.5) node[above] {Voltage $U$};
		
		% Straight line
		\draw[thick] (0,0) -- (5,3.5);
		
		% Measurement points
		\fill (1,0.7) circle (2pt);
		\fill (2,1.4) circle (2pt);
		\fill (3,2.1) circle (2pt);
		\fill (4,2.8) circle (2pt);
		\fill (5,3.5) circle (2pt);
		
		% Slope triangle
		\draw[dashed] (2,1.4) -- (4,1.4);
		\draw[dashed] (4,1.4) -- (4,2.8);
		
		% Label
		\node[right] at (4,2.1) {$R=\frac{\Delta U}{\Delta I}$};
		
	\end{tikzpicture}
	\caption{The proportionality between voltage and current becomes visible from the measurement points.
		The slope of the straight line corresponds to the resistance $R$.}
	\label{fig:ohm_proportionality}
\end{figure}

Other experiments also reveal such structures.
In a thermometer, the approximately uniform expansion of a liquid
leads to an approximately linear relationship.
In the work of \textbf{Galileo Galilei}, a different form appears instead:
motion on an inclined plane and projectile motion lead to parabolic relationships.
Thus the experiment shows whether nature follows a linear,
quadratic, or other mathematical pattern.

This makes a fundamental idea of natural science visible:
thinking alone does not decide which mathematics fits the world.
Only the experiment shows
which structure is actually present in reality.

\begin{NoteBox}[Experiment as Translator]
	\label{box:experiment}
	\small
	The experiment connects measurement and mathematics:
	it shows whether a natural phenomenon follows a linear, quadratic,
	exponential, or other structure.
	Only in this way does an abstract formula become a physical law.
\end{NoteBox}

% =====================================================
% 1.4 Conclusion
% =====================================================
\subsection{Conclusion}

At the beginning of mathematics stand the numbers.
Already with them begins the path from immediate experience
to abstract structure.
From simple counting through negative and irrational numbers
all the way to complex numbers, one thing becomes clear:
each expansion of the concept of number opens a new access
to the order of the world.

This condenses the guiding idea of this book.
Mathematics is not merely a tool invented by human beings,
but appears as a structure
that is discovered again and again
under different cultural and historical conditions.
Its validity does not depend on language,
place, or time,
but on the inner necessity of its concepts and relationships.

For this reason, this book begins with numbers.
They are not only the first object of mathematics,
but its most original access to reality.
Chapter II begins precisely here
and follows the fundamental building blocks
from which the further mathematical structure unfolds.

\begin{NoteBox}[Numbers as a Universal Discovery]
	\label{box:universal}
	\small
	The history of mathematics shows:
	different cultures independently arrived
	at the same fundamental ideas of number and calculation.
	The path into mathematics therefore does not begin with a mere convention,
	but with a structure that is discovered again and again.
\end{NoteBox}